\documentclass[pdflatex,sn-nature]{sn-jnl}

\usepackage{graphicx}%
\usepackage{multirow}%
\usepackage{amsmath,amssymb,amsfonts}%
\usepackage{amsthm}%
\usepackage{mathrsfs}%
\usepackage[title]{appendix}%
\usepackage{xcolor}%
\usepackage{lineno}%
\usepackage{textcomp}%
\usepackage{manyfoot}%
\usepackage{booktabs}%
\usepackage{algorithm}%
\usepackage{algorithmicx}%
\usepackage{algpseudocode}%
\usepackage{listings}%

\begin{document}

\title[Southern Ocean WMT Across Climate States]{Glacial–interglacial regime shift in Southern Ocean dense water formation}


%%------------------------------------------------------------
%%  Authors and affiliations
%%------------------------------------------------------------
\author*[1]{\fnm{Xiaoxu} \sur{Shi}}\email{Xiaoxu Shi (shixiaoxu@sml-zhuhai.cn)}
\author*[2,1]{\fnm{Jiping} \sur{Liu}}\email{Jiping Liu (liujp63@mail.sysu.edu.cn)}
%\author*[2,1]{\fnm{Jiping} \sur{Liu}}\email{Jiping Liu (liujp63@awi.de)}
\author[1]{\fnm{Chaoyuan} \sur{Yang}}
\author[3]{\fnm{Gerrit} \sur{Lohmann}}

%\affil[1]{\orgname{Alfred Wegener Institute, Helmholtz Center for Polar and Marine Research}, \orgaddress{\city{Bremerhaven}, \country{Germany}}}
\affil[1]{\orgname{Southern Marine Science and Engineering Guangdong Laboratory (Zhuhai)}, \orgaddress{\city{Zhuhai}, \country{China}}}
\affil[2]{\orgname{School of Atmospheric Sciences, Sun Yat-sen University}, \orgaddress{\city{Zhuhai}, \country{China}}}
\affil[3]{\orgname{Alfred Wegener Institute, Helmholtz Center for Polar and Marine Research}, \orgaddress{\city{Bremerhaven}, \country{Germany}}}

\linenumbers


\abstract{Winter dense water formation in the Southern Ocean drives global overturning and regulates deep-ocean heat and carbon storage, but how surface buoyancy forcing controls this process across past climate states remains unclear. Using water mass transformation (WMT) analysis in AWI-ESM simulations of the pre-industrial, mid-Holocene, Last Interglacial, Last Glacial Maximum, and Marine Isotope Stage 3, we identify a sea-ice-mediated regime shift between climate states. In the model, interglacial dense water formation is driven mainly by turbulent heat loss over the open ocean, whereas glacial formation is haline-driven by brine rejection as expanded sea ice insulates the ocean from the atmosphere. Further analysis suggests that the Southern Annular Mode modulates transformation by 3--8 Sv between its positive and negative phases in every climate state, working through heat loss when the ocean is open and through brine rejection when it is ice covered. Weaker glacial WMT coincides with a collapse of the simulated abyssal overturning cell from about 11 to 2 Sv and with older abyssal water. Ideal ages at 4000 m remain below 500 years during interglacials but exceed 750 years basin-wide during glacials, reaching 1500 years in the Pacific. Our results provide a mechanistic framework for interpreting past Southern Ocean circulation, while also showing that the simulated pre-industrial state transforms water at lighter densities than observations indicate.}

%\abstract{Winter dense water formation in the Southern Ocean drives global overturning circulation and regulates deep-ocean heat and carbon storage. However, understanding how surface buoyancy forcing affects this formation across different climate states remains a fundamental challenge. Using water mass transformation (WMT) analysis in AWI-ESM simulations spanning pre-industrial, mid-Holocene, last interglacial, last glacial maximum, and marine isotope stage 3, we identify a fundamental regime shift between climate states.  Interglacial WMT is thermally dominated, driven by surface heat loss, whereas glacial WMT is haline-dominated, governed by expanded sea ice through brine rejection. Superimposed on these mean states, the Southern Annular Mode (SAM) systematically modulates wintertime WMT across all  climate states, with anomalies of 3-8 Sv between high- and low-SAM winters. During positive SAM phases, dense water formation is intensified via enhanced heat loss over open ocean and amplified brine rejection within coastal areas. Thermal forcing dominates the SAM-induced response during interglacials, while haline forcing dominates during glacials.  In addition, simulated deep-ocean ventilation ages remain below 500 years in most basins during interglacials but exceeding 750 years basin-wide during glacials, with Pacific ages reaching 1500 years. These results highlight the state-dependence of surface forcing, offering a mechanistic framework for understanding past Southern Ocean circulation changes.}


%Here, we analyze water mass transformation (WMT) in the AWI-ESM coupled model across five climate scenarios spanning pre-industrial (PI), mid-Holocene (MH), Last Interglacial (LIG), Last Glacial Maximum (LGM), and Marine Isotope Stage 3 (MIS3). By systematically decomposing  surface buoyancy fluxes into thermal (radiative and turbulent heat) and haline (sea ice, net precipitation, and runoff) components,  we identify a fundamental regime shift in forcing mechanisms. Specifically, interglacial WMT is thermally dominated, with surface heat loss being the primary contributor of dense water formation. Conversely, glacial states exhibit haline-dominated forcing, where expanded sea ice governs WMT through combined effects of enhanced heat loss and brine rejection. Superimposed on these mean states, we further demonstrate that the Southern Annular Mode (SAM) systematically modulates wintertime WMT across all climate states, with maximum anomalies of 3-8 Sv between high and low SAM winters.  During positive SAM phases, dense water formation is intensified via enhanced heat loss over open ocean and amplified brine rejection within coastal areas. Thermal forcing dominates the SAM-induced response during interglacials, while haline forcing dominates during glacials.  Reflecting the shift in ventilation efficiency, simulated deep ocean ventilation ages increase substantially during glacial periods, with Pacific ages reaching 1500 years compared to less than 500 years during interglacials with marked regional heterogeneity, consistent with radiocarbon proxy reconstructions. These results highlight the non-linear state-dependence of surface forcing, offering a mechanistic framework for interpreting past Southern Ocean circulation changes.

\keywords{Southern Ocean, water mass transformation, paleoclimate, sea ice, Southern Annular Mode, glacial-interglacial}

\maketitle

% ============================================================================
% Introduction
% ============================================================================
\section{Introduction}

%\subsection{The Southern Ocean and global climate}

The Southern Ocean plays an important role in Earth's climate system. It serves as the primary region where deep waters return to the surface through wind-driven upwelling and where surface waters are transformed into dense water masses that ventilate the global abyssal ocean \cite{Marshall2012,Talley2013}. This region is characterized by intense air-sea interaction, extensive seasonal sea ice coverage, and strong meridional density gradients that drive both the upper and lower cells of the global meridional overturning circulation (MOC) \cite{Speer2000,Lumpkin2007}. Through this overturning circulation, the Southern Ocean regulates the global heat and carbon storage \cite{Marinov2006,Gray2024}. The efficiency of this regulation depends critically on the rate of dense water formation and its subsequent isolation from atmospheric contact. Therefore, understanding physical processes behind dense water formation in the Southern Ocean across different climate states is essential for interpreting past and future ocean circulation changes and their climate impacts.

%Quantitatively, the Southern Ocean accounts for approximately 40\% of global anthropogenic CO${_2}$ uptake \cite{Gray2024}, underscoring its central role in modulating atmospheric CO${_2}$ concentrations. Therefore, understanding the physical processes behind dense water formation in the Southern Ocean is crucial for predicting how the climate system will respond to future changes. %Through this overturning circulation, the Southern Ocean regulates the global capacity to store heat and carbon \cite{Marinov2006,Gray2024}.  As the primary pathway for deep water ventilation, it mediates carbon exchange between the abyssal reservoir and the atmosphere, with sequestration efficiency depending critically on the rate of dense water formation and its subsequent isolation from atmospheric contact. 


Surface buoyancy fluxes fundamentally control the transformation of water masses in density space \cite{Walin1982,Speer1992,drake2025water}. These fluxes consist of heat exchange with the atmosphere and freshwater fluxes associated with precipitation, evaporation, river runoff, and sea ice processes. In the Southern Ocean, the surface buoyancy fluxes exhibit pronounced spatial and seasonal variability. Specifically, intense wintertime cooling drives surface densification in open water areas, whereas sea ice formation and melting produce strong haline forcing through brine rejection and freshwater release \cite{Abernathey2016,Pellichero2018}. Observational evidence suggests that freshwater fluxes, particularly those associated with sea ice thermodynamics, may dominate in ice-covered regions \cite{Pellichero2018,Bailey2023}. However, quantifying the relative contribution of thermal versus haline forcing to dense water formation across the entire Southern Ocean remains unclear.

%The role of the Southern Ocean in the global carbon cycle is equally important. As the primary pathway for deep water ventilation, this region regulates carbon exchange between the deep ocean reservoir and the atmosphere \cite{Marinov2006,Gray2024}. The efficiency of this carbon sequestration process depends critically on the rate at which surface waters are transformed into dense bottom waters and subsequently isolated from direct atmospheric contact. Quantitatively, the Southern Ocean accounts for approximately 40\% of the global ocean's anthropogenic CO$_2$ uptake \cite{Gray2024}, underscoring its central role in modulating atmospheric CO$_2$ concentrations and climate change.


%\subsection{The water mass transformation framework}

The water mass transformation (WMT) framework provides a thermodynamic approach to quantifying how air-sea fluxes drive changes in water mass properties \cite{Walin1982,Groeskamp2019,drake2025water}. This method relates surface heat fluxes to cross-isothermal mass transport \cite{Walin1982}, and was subsequently extended  to incorporate both thermal and freshwater contributions to surface density flux \cite{Speer1992}.

WMT implementations for the modern climate decompose surface fluxes into multiple components to isolate specific physical processes \cite{Iudicone2008a,Iudicone2008b}. Heat flux contributions are separated into radiative (shortwave and longwave) and turbulent (sensible and latent heat) components. Freshwater flux decomposition distinguishes between sea ice formation/melt, glacial meltwater, and meteoric sources such as precipitation and runoff \cite{Abernathey2016,Bailey2023}. The sea ice component is particularly important, as it acts as a freshwater pump that removes freshwater from the Southern Ocean through equatorward ice transport and melting, concentrating salt in formation regions \cite{Abernathey2016}.

Observation-based WMT analyses have improved our understanding of Southern Ocean water mass formation.  Based on Argo floats, ship data, and instrumented marine mammals, a previous study showed that seasonal sea ice growth and melt dominates transformation in the ice-covered regions, with freshwater fluxes drive the meridional overturning that feeds dense water formation \cite{Pellichero2018}. Additionally, the closed-budget WMT analysis in the Weddell Sea reveals that sea ice brine rejection dominates the surface salt flux for most seasons \cite{Bailey2023}.


The paleoclimate record, together with model simulations, offers a means for understanding Southern Ocean responses to climate forcings beyond the range of modern observations. Across glacial-interglacial cycles, atmospheric CO$_2$ \cite{SigmanBoyle2000,Kohfeld2005}, ice sheet extent \cite{huybrechts1990antarctic}, and deep ocean circulation \cite{rafter2022global} all varied dramatically. Radiocarbon-based reconstructions indicate that these circulation changes substantially altered deep water ventilation, with ventilation age increasing by $\sim$689$\pm$53 $^{14}$C-yr during glacial periods \cite{Skinner2017}, a slowdown closely tied to deep ocean carbon sequestration that may account for more than half of the glacial-interglacial CO$_2$ change \cite{Skinner2017}. A consistent mechanistic explanation for these changes is emerging from theory and modeling. Glacial Antarctic sea ice expansion is expected to enhance the haline contribution to surface buoyancy forcing through increased brine rejection, while simultaneously insulating the deep ocean from atmospheric forcing, thereby altering both the rate and properties of dense water formation \cite{Ferrari2014}. Coupled glacial simulations confirm that this brine-rejection-driven surface buoyancy loss exerts a strong control on abyssal stratification and overturning, although its magnitude varies substantially between models depending on how sea-ice formation is represented \cite{Marzocchi2017}. This points to a thermal-to-haline shift in surface forcing as climate cools, which needs to be quantified within a consistent WMT framework across a range of climate states. Indeed, existing paleoclimate modeling studies have largely diagnosed circulation changes through overturning streamfunctions, water mass volumes, or deep-ocean tracer distributions \cite{Weber2007,Wainer2012}. The surface buoyancy fluxes that drive those changes have rarely been decomposed explicitly. A recent study using the same ocean model in a forced, ocean-only configuration attributes the large glacial volume of Atlantic bottom water to enhanced sea ice export and to weaker mixing between northern- and southern-sourced deep water \cite{Chen2025}. That result identifies sea ice as the central agent from the distribution of water masses and from the overturning streamfunction. It does not measure the surface transformation rates themselves, which is what a decomposed buoyancy budget provides and what allows the thermal and haline pathways to be compared directly and across several climate states.

Beyond the mean state, dense water formation also responds to internal atmospheric variability. The Southern Annular Mode (SAM) is the leading mode of extratropical Southern Hemisphere circulation variability, describing a meridional redistribution of atmospheric mass between mid-latitudes and Antarctica \cite{Gong1999,Thompson2000,Marshall2003}. Its positive phase strengthens the circumpolar westerlies and displaces them poleward, which changes Ekman divergence, air-sea heat exchange, and the northward export of sea ice \cite{HallVisbeck2002,SenGupta2006}. In the present-day ocean these adjustments modulate mixed-layer depth and surface buoyancy forcing \cite{Sallee2010,SenGupta2006}. Surface buoyancy forcing in turn sets the rate at which surface waters are transformed into denser classes \cite{Cerovecki2013,Abernathey2016}. The SAM has also been linked to observed changes in bottom water properties \cite{Zhou2023,Silvano2020}. Whether the SAM remains an effective modulator of dense water formation once the mean state changes, in particular once an expanded ice cover separates the ocean from the atmosphere, is not known.

%The paleoclimate record and model simulations provide essential context for understanding Southern Ocean responses to climate forcings beyond the range of modern observations. Glacial-interglacial cycles feature dramatic changes in atmospheric CO$_2$ \cite{SigmanBoyle2000,Kohfeld2005}, ice sheet extent \cite{huybrechts1990antarctic}, and deep ocean circulation \cite{rafter2022global}. Proxy evidence indicates substantially altered deep water ventilation rates, i.e., radiocarbon-based reconstructions suggested that ventilation age increases of  $\sim$689$\pm$53 ${14}$C-yr during glacial periods \cite{Skinner2017}. These circulation changes are linked to the efficiency of deep water formation and ocean carbon sequestration, potentially explaining more than half of the glacial–interglacial CO$_2$ changes \cite{Skinner2017}.  Existing paleoclimate modeling studies have largely diagnosed circulation changes through overturning streamfunctions, water mass volumes, or deep-ocean tracer distributions \cite{Weber2007,Wainer2012}, without explicitly decomposing the surface buoyancy fluxes that drive those changes. Theoretical work further suggested that glacial Antarctic sea ice expansion enhances haline contributions to surface buoyancy forcing via increased brine rejection  while insulates deep waters from atmospheric forcing, thereby altering the rate and properties of dense water formation \cite{Ferrari2014}.  Coupled glacial simulations confirm that Antarctic sea-ice formation, through brine-rejection-driven surface buoyancy loss, exerts a significant control on abyssal stratification and overturning, but the magnitude of this effect varies substantially across models depending on their representation of sea-ice formation \cite{Marzocchi2017}. This regime transition has not been quantified within a consistent WMT framework across a range of climate states. Furthermore, while the Southern Annular Mode (SAM) is known to modulate present-day Southern Ocean mixed-layer depth and air-sea heat exchange\cite{Sallee2010}, whether this atmospheric pattern remains a robust modulator of dense water formation under different climate boundary conditions has not been examined. Bridging this gap requires explicit WMT decomposition across multiple climate states, spanning interglacials to glacials, to isolate the individual forcing components and assess the state-dependence of both mean-state and internally driven variability.

%Despite recent progress, systematic WMT application to paleoclimate simulations remains limited.  

%Despite current progress, systematic WMT analysis across multiple paleoclimate states remains limited. Most paleoclimate modeling studies have focused on circulation strength or water mass properties rather than the underlying surface forcing mechanisms \cite{Weber2007,shi2020early}. Understanding how surface buoyancy fluxes drive dense water formation and how this mechanism changes across climate states requires explicit WMT decomposition that isolates individual forcing components.

%Ferrari2014: the expansion of permanent sea ice in the Southern Hemisphere results in a volume increase of Antarctic-origin abyssal waters and a reduction in mixing between abyssal waters of Arctic and Antarctic origin.


%The paleoclimate perspective is particularly valuable as the modern instrumental record provides insufficient temporal scope to evaluate how ocean processes respond to large forcing changes \cite{Tierney2020Paleo}. Paleoclimate states spanning different CO$_2$ levels, ice sheet configurations, and sea ice extents provide out-of-sample tests for understanding process sensitivities. Specifically, for Southern Ocean dense water formation,  paleoclimate analysis offers the only means to evaluate how the balance of thermal versus haline forcing shifts under different boundary conditions.


This study analyzes WMT in five AWI-ESM simulations spanning a range of interglacial and glacial climate states, by decomposing surface buoyancy fluxes into radiative, turbulent, and haline components to identify the mechanisms governing dense water formation across interglacial and glacial regimes. We further examine how the Southern Annular Mode  modulates WMT rates and whether its coupling with surface buoyancy forcing remains robust across varying climate states. Finally, deep ocean ventilation ages are assessed as an integral diagnostic of overturning efficiency. This analysis serves as  an integral proxy for evaluating the efficiency of the overturning circulation in response to surface forcing shifts.

%Finally, to link these surface processes to large-scale circulation changes, we assess deep ocean ventilation ages and their regional heterogeneity. This diagnostic serves as  an integral proxy for evaluating the efficiency of the overturning circulation in response to surface forcing shifts. By focusing on the mechanistic drivers of buoyancy gain  and loss rather than specific water mass definitions, this framework provides robust physical constraints on how Southern Ocean ventilation responds to past, and potentially future, climate perturbations.

%In section 2, we introduce the model used in our study and the methods for surface buoyancy flux decomposition and water mass transformation. Results of


\section{Results}

\subsection{Large-scale features of the Southern Ocean during austral winter}

The figures label the five simulated climate states as PI (pre-industrial), MH (mid-Holocene), LIG (Last Interglacial), LGM (Last Glacial Maximum) and MIS3 (Marine Isotope Stage 3). The text uses the full names for the mid-Holocene and the Last Interglacial. Their boundary conditions are summarised in Table \ref{tab:bc}.


To understand the surface preconditioning for dense water formation, we first examine the surface thermodynamic and dynamic state during austral winter, the primary season for this process in the Southern Ocean. The maps in Fig. \ref{sst} are averaged over June to September. % The simulated JJA climate patterns reveal distinct characteristics across the five simulated periods (Fig. \ref{sst}-\ref{wind}).
The simulated surface anomalies relative to PI (Fig. \ref{sst}) show a striking contrast between interglacial (mid-Holocene, Last Interglacial) and glacial (LGM, MIS3) states. During the interglacial periods, anomalies are generally moderate. The mid-Holocene shows weak SST and SSS changes across the Southern Ocean (Fig. \ref{sst}a, e). In the Last Interglacial, the most important feature is a distinct warm-salty pattern in the Bellingshausen-Amundsen area, where pronounced warming (more than 2 °C) is density-compensated by concurrent salinification (Fig. \ref{sst}b, f). This compensation results in localized surface densification (Fig. \ref{wind}b).
In  contrast, the glacial periods (LGM, MIS3) are characterized by profound basin-wide cooling (Fig. \ref{sst}c,d) and salinification (Fig. \ref{sst}g,h). Here the enhanced glacial salinification partly reflects our model initialization with higher global mean salinity, accounting for ice sheet storage of freshwater. Because that offset is applied uniformly, the anomaly maps in Fig. \ref{sst}g,h combine it with the locally generated signal, so only their regional structure reflects local dynamics. The combination of extreme cooling and salinification leads to a pronounced increase in surface density across the entire Southern Ocean (Fig. \ref{wind}c,d), particularly in the Pacific sector. In the glacial states the surface density anomalies (Fig. \ref{sst}i-l) follow the salinity field. Cooling and salinification both act to densify the surface, but at the near-freezing temperatures that prevail south of the ice edge the thermal expansion coefficient is small, so the haline term sets both the pattern and the magnitude. In the Last Interglacial the warm and salty Bellingshausen-Amundsen anomaly is close to density compensated (Fig. \ref{sst}j), and the LGM retains weak negative density anomalies in parts of the Atlantic-Indian sector where surface freshening outweighs cooling.

The sea ice anomalies (Fig. \ref{sst}i-l) follow the temperature changes. In the mid-Holocene the changes are small. In the Last Interglacial the winter sea ice decreases strongly in the Ross sector, where the winter ice edge retreats south by about 5° of latitude beneath the warm and salty surface anomaly. In LGM and MIS3 the winter sea ice area expands substantially and the ice edge moves north by about 10--15° of latitude, most of all in the Weddell sector.


In the interglacial states (PI, mid-Holocene, Last Interglacial), the model produces deep mixed layers in the Weddell and Ross Sea gyres (Fig. \ref{mld}a-c), indicating active open-ocean convection. In PI this convection is concentrated in the Weddell sector. Conversely, under glacial conditions (LGM, MIS3), the glacial cooling drives a dramatic northward expansion of sea ice extent (Fig. \ref{mld}i, j), covering the majority of the subpolar Southern Ocean (Fig. \ref{mld}i-j). This massive sea ice expansion creates an insulation effect, preventing the water column from direct heat loss to the atmosphere. As a result, the mixed layer depth is strongly reduced and becomes confined to narrow coastal bands and polynyas (Fig. \ref{reso}). The reduction is substantial, and some deep mixed layers remain. Summed south of 55°S, the area with a winter mixed layer deeper than 400 m falls from $7.5\times10^{11}$~m$^2$ in PI to $1.7\times10^{11}$~m$^2$ in LGM and $1.6\times10^{11}$~m$^2$ in MIS3, and the maximum mixed-layer depth decreases from 991 m to 819 m and 645 m. Mixed layers deeper than 600 m, which cover $2.4\times10^{11}$~m$^2$ in PI, essentially disappear in both glacial states. The residual glacial signals in coastal regions are therefore genuinely shallower than their interglacial counterparts (Supplementary Fig. \ref{s1}), consistent with dense water production that is confined to narrow coastal bands. %This spatial reorganization of MLD suggests a fundamental regime shift in the mechanism of dense water formation: transitioning from widespread open-ocean cooling in interglacials to processes dominated by brine rejection in coastal polynyas during glacials.


These thermodynamic changes are accompanied by shifts in atmospheric forcing (Fig. \ref{sst}m-p). The mid-Holocene shows moderate anomalies in zonal wind, generally negative south of about 50°S and positive to the north. The Last Interglacial shows a similar pattern but with a strong intensification in the Pacific sector (Fig. \ref{sst}n). Both LGM and MIS3 exhibit strong positive anomalies in zonal wind speed and wind stress ($>0.04$~N/m$^2$) across the circumpolar belt, although the LGM also shows some negative anomalies in the Atlantic-Indian sector (Fig. \ref{sst}o,p; Fig. \ref{wind}g,h). This intensified wind stress enhances Ekman transport and facilitates the northward expansion of sea ice, fundamentally altering the air-sea interaction interface. In the glacial states the sea ice increases all around the continent. Averaged between 60 and 75°S, winter ice concentration increases in the Ross, Weddell and Ad\'elie sectors at the Last Glacial Maximum, and all three sectors cool to within about 0.1~K of the surface freezing point, so ice growth is no longer limited by the surface heat content anywhere. In the Last Interglacial the Ross sector has the most pronounced change. There the surface ocean between 60 and 75°S is warmer than in PI and sits well above the freezing point, and winter ice concentration falls. In that state the wind can export ice efficiently but the water is too far from freezing for the exported ice to be replaced, so the sector loses ice cover and gains exposure to the atmosphere. This asymmetry is why the Ross sector departs from the circumpolar pattern in the variability analysis below.

The simulated westerlies strengthen in both the glacial and the Last Interglacial states, for different reasons. Measured as the latitude and magnitude of the maximum winter zonal-mean 10 m zonal wind, the wind maximum sits at 45.7°S with 5.83~m~s$^{-1}$ in PI and is unchanged in the mid-Holocene. It shifts to 47.6°S and 6.92~m~s$^{-1}$ in the Last Interglacial and to 49.4°S with 6.30 and 6.85~m~s$^{-1}$ in LGM and MIS3. In the glacial states the change follows the expanded sea ice. The winter ice edge moves from 66°S in PI to 55--57°S, and the strongest near-surface temperature gradient moves with it from about 65°S to 60--62°S and becomes 24--29\% steeper. This gradient still lies 10--12° south of the wind maximum. The lower atmosphere therefore becomes more baroclinic on the poleward side of the westerlies, and the Eady growth rate between 850 and 500~hPa increases by 21--28\% between 45 and 65°S. The tropical upper troposphere also cools more than the surface, which weakens the temperature gradient aloft in mid-latitudes. Both changes favour stronger westerlies further south. In the Last Interglacial the change comes from the orbital forcing. Winter insolation increases by 26--33~W~m$^{-2}$ between 30 and 45°S, the tropical upper troposphere warms by about 2~K, and the temperature gradient aloft steepens, while the circumpolar mean sea ice edge retreats by about 1°. The westerlies therefore strengthen with a small poleward shift. Five equilibrium states cannot separate the roles of sea ice, ice sheets and greenhouse gases, which would require atmosphere-only sensitivity experiments. That the mid-Holocene remains indistinguishable from PI argues against residual model drift as the origin of these changes.


\subsection{Water mass transformation analysis}


The decomposition of WMT rates reveals a fundamental regime shift in dense water formation mechanisms between interglacial and glacial climate states (Fig. \ref{wmt_jja}).

In our model, dense water formation is predominantly thermally driven in the interglacial experiments. In the austral winter (JJA), when transformation rates peak, the Southern Ocean exhibits a WMT maximum of 70--81~Sv at $\sigma_2 \approx 36.4$-$36.8$ kg m$^{-3}$ across all simulated interglacial states (Fig. \ref{wmt_jja}). This transformation is dominated by surface heat loss (red line), which accounts for 70--85\% of the total at the transformation maximum, with limited sea ice extent restricting the haline contribution from brine rejection to a secondary role (15--30\%; green line). Other freshwater fluxes (blue lines) oppose densification through net precipitation and river runoff. We note that this thermally dominated maximum lies at the densities of intermediate and mode waters, which are lighter than bottom water. In the densest classes that the simulated surface fluxes reach ($\sigma_2 \ge 37.0$~kg~m$^{-3}$), the PI transformation is already haline dominated, with the sea ice term contributing 0.79~Sv against 0.06~Sv from heat in the annual mean. The thermal dominance reported here is therefore a property of the integral over all density classes, and we quantify its sensitivity to the choice of integration domain and density range in the Discussion.
% claude（11-27，0.1 格）：本段 70--85\%、15--30\% 和 0.79/0.06 Sv 没改，待你定（见 pass_20261127c/bin0p1_table.md）。

Glacial periods exhibit a distinctly different regime. The winter WMT maximum weakens substantially to 50-59 Sv while shifting to higher density classes ($\sigma_2 = 37.1$-$37.6$ kg m$^{-3}$; Fig. \ref{wmt_jja}). More fundamentally, the dominant forcing mechanism reverses, with sea ice processes contributing almost the entire of  the total WMT. This haline dominance reflects expanded sea ice coverage, which exerts dual effects on dense water formation.  Extensive ice cover insulates the ocean surface from atmospheric heat exchange, suppressing thermally driven transformation. On the other hand, enhanced sea ice production concentrates brine rejection, intensifying haline-driven densification at higher density classes ($\sigma_2 > 37$ kg m$^{-3}$). Thermal forcing operates only secondarily.

This interglacial thermal-dominance versus glacial haline-dominance signature is consistent across regional sectors, including the Ross Sea, Weddell Sea, and Ad\'{e}lie regions. Note that the vertical axes in Fig. \ref{wmt_jja} are scaled independently for each panel to make the shape of the transformation curves legible; a version with a common vertical axis, which allows the magnitudes to be compared directly between climate states, is provided as Supplementary Fig. \ref{s_common_axis}.

%This thermally dominated regime persists in the annual mean (Fig.~\ref{wmt_ann}), where net heat loss maintains a transformation rate of up to 18 Sv, balancing the seasonal buoyancy gain from summer warming.
%The glacial states mark a transition to a haline-modulated regime. While intense winter cooling remains the primary driver of instantaneous buoyancy loss (Fig.~\ref{wmt_jja}), the resulting water masses are shifted to significantly denser classes ($\sigma_2 > 37.0$~kg m${-3}$). Crucially, the role of sea ice freshwater flux shifts from a background player to a leading driver. In the annual mean (Fig.~\ref{wmt_ann}), the sea ice contribution (green line) in the LGM and MIS3 scenarios dramatically intensifies, becoming the dominant positive forcing term in the Southern Ocean and Weddell Sea sectors. This indicates that while heat loss drives the seasonal convection cycle, the expanded sea ice engine---via enhanced brine rejection---is the critical mechanism that locks these waters into the abyssal density classes.


%Regional Heterogeneity
%This mechanistic shift exhibits strong spatial heterogeneity. The Weddell Sea, serving as the classical engine of deep water formation, shows the most pronounced glacial intensification. Here, winter transformation increases to $\sim$30~Sv (at $\sigma_2 = 37.2$~kg m${-3}$), heavily reinforced by a threefold increase in sea ice-driven transformation compared to the PI state. In contrast, the Ross Sea maintains a quasi-thermal balance even under glacial conditions, suggesting that local polynya dynamics allow strong atmosphere-ocean heat fluxes to persist despite the broader expansion of the sea ice edge.


\subsection{Surface density tendency from buoyancy forcing}

To identify the geographic origins and specific physical components driving the identified WMT regimes, we convert heat and freshwater fluxes into surface density tendency during austral winter. We express it as a surface density flux in kg~m$^{-2}$~month$^{-1}$, which equals the density change of a surface layer 1 m thick over one month (see Methods).

In PI, strong thermal densification (red shading) is concentrated in the open ocean (Fig. \ref{density_heat}a), with maximal values reaching about 13~kg~m$^{-2}$~month$^{-1}$ in the circumpolar band between 50-65°S. This thermal densification is overwhelmingly controlled by turbulent fluxes (sensible and latent heat), with radiative fluxes playing a minor role in the buoyancy budget (Fig. \ref{density_heat}f-j vs. Fig. \ref{density_heat}k-o), consistent with observationally based surface buoyancy budgets for the Southern Ocean \cite{Cerovecki2013,Abernathey2016}.
Interglacial anomalies (mid-Holocene, Last Interglacial; Fig. \ref{density_heat}b,c,g,h,l,m) show small to modest changes relative to PI north of the ice-covered region all around Antarctica. The mid-Holocene exhibits no obvious thermal density tendency anomalies in the ice-covered region, while
The Last Interglacial exhibits pronounced increases in thermal-driven surface density tendency  localized to the Ross Sea sector. This arises from the local ice-free condition during the Last Interglacial, which facilitates enhanced ocean-atmosphere heat flux in the Ross Sea, a region that remains ice-covered throughout PI and mid-Holocene winter months.
For glacial periods (LGM, MIS3; Fig. \ref{density_heat}d,e), despite colder atmospheric temperatures, thermal density tendency decreases throughout the Southern Ocean. Anomalies reach about $-5$~kg~m$^{-2}$~month$^{-1}$ in the Pacific sector. This result arises from expanded glacial sea ice acting as a thermal insulating lid that prevents ocean-atmosphere heat exchange. The turbulent flux component (Fig. \ref{density_heat}n,o) shows the strongest reduction, confirming that ice cover blocks sensible and latent heat loss more effectively than radiative exchange.

The sea ice contribution shows a dipole in space. Fig. \ref{density_fwf}f exhibits this classic pattern in the sea ice-driven density tendency.  Coastal divergence zones experience strong brine-driven densification while ice edge areas show weak buoyancy gain from ice melt. This pattern dominate the total haline density tendency (Fig. \ref{density_fwf}a). Other freshwater sources (net precipitation and runoff, Fig. \ref{density_fwf}k) contribute negligibly to the surface density tendency in these simulations. In the present-day ocean, meltwater derived from the Antarctic ice sheet, both as basal melt beneath ice shelves and as iceberg discharge, is an important part of the coastal freshwater budget and acts to suppress dense water formation \cite{Silvano2018,ToggweilerSamuels1995}. Because the model has no ice shelf cavities and no prescribed meltwater flux, basal melt and iceberg melt are not represented as such. The freshwater that leaves the ice sheet returns to the ocean as liquid runoff at the coast (Methods).
Mid-Holocene anomalies remain small (Fig. \ref{density_fwf}b,g,l), consistent with limited sea ice changes during the mid-Holocene compared to PI. During the Last Interglacial, the reduced sea ice in the Ross Sea contributes to a weakening of brine rejection and therefore a weakened density tendency (Fig. \ref{density_fwf}c,h).
Glacial periods show a pronounced northward expansion and intensification of the buoyancy loss signatures relative to PI (Fig. \ref{density_fwf}d,e,i,j). Total freshwater density tendency anomalies reach about $+6.5$~kg~m$^{-2}$~month$^{-1}$ in a circumpolar band between 55 and 70°S. The sea ice component (Fig. \ref{density_fwf}i,j) dominates this signal, with the strongest anomalies along the coast, where they exceed 10~kg~m$^{-2}$~month$^{-1}$ over a tenth of the area south of 65°S and locally reach about 47~kg~m$^{-2}$~month$^{-1}$.


\subsection{Southern Annular Mode impact on water mass transformation}


We now examine how the SAM, introduced above, modulates dense water formation across different climate states and regional sectors. We compare high-SAM and low-SAM winters in the 100-year monthly simulations of each state. High-SAM (low-SAM) years are those in which the JJA-mean SAM index lies more than a fixed multiple of its standard deviation above (below) its climatological mean, and the threshold used for each composite is given in the figure caption. The supporting composite fields are collected in Supplementary Figs. \ref{sam_wmt}--\ref{sam_windice}.

Across all climate states, the positive phase of SAM systematically enhances WMT in the Southern Ocean (Fig.~\ref{sam_wmt}a-e, black lines), indicating that a poleward shift and intensification of the westerlies favors dense water production. However, the magnitude and mechanism of this SAM-WMT coupling differ substantially between interglacial and glacial periods. For PI, the maximum impacts occur at $\sigma_2 = 36.7$ kg/m$^3$ with the SAM-induced WMT anomaly being 3.2 Sv,  while in the mid-Holocene and the Last Interglacial, the maximum WMT anomaly between high and low SAM years can reach about 8 Sv, occurring at $\sigma_2 = 36.5$ kg/m$^3$, representing the strongest SAM-WMT coupling among all climate states. During glacial periods, peak anomalies are more modest (+4 Sv for LGM at $\sigma_2 = 37.5$~kg/m$^3$; +4.2 Sv for MIS3 at $\sigma_2 = 37.1$~kg/m$^3$) and shifted  toward higher density classes, reflecting colder and saltier background conditions.
% claude（11-27，0.1 格）：本段 SAM 数字（3.2、8、4、4.2 Sv）连 0.2 格都对不上；0.1 格最大异常 PI 6.1、MH 8.2、LIG 10.4、LGM 10.2、MIS3 7.0 Sv，冰期不再「更小」，待你定。摘要和讨论的 3--8 Sv 同。


The partitioning between thermal and haline contributions reveals a fundamental  climate-state dependence in the SAM-WMT relationship. During interglacials, the thermal component (red lines) dominates the WMT response while the sea ice freshwater forcing plays a secondary role, especially in the Ross Sea and Weddell Sea sectors (Fig.~\ref{sam_wmt}a-c,f-h,k-m). This thermal dominance reflects the mechanism whereby strengthened westerlies during positive SAM phases enhance turbulent heat loss at the air-sea interface, driving densification through surface cooling. In contrast, during glacial periods, sea ice freshwater flux (green lines) emerges as the primary contributor (Fig.~\ref{sam_wmt}d,e,i,j,n,o,s,t). This indicates that under extensive glacial sea ice cover, SAM-driven changes in brine rejection become the dominant process for modulating dense water formation.

Regional decomposition reveals that the Weddell Sea and Ad\'elie sectors contribute consistently to the circumpolar SAM-WMT signal across interglacial climate states (Fig.~\ref{sam_wmt}k-m,p-r), while the Ross Sea exhibits more variable and sometimes opposing responses (Fig.~\ref{sam_wmt}f-h). Moreover, the glacial WMT response to SAM is overall dominated by the Weddell Sea region (Fig.~\ref{sam_wmt}n,o), whereas the Ross Sea and Ad\'elie sectors play a minor role (Fig.~\ref{sam_wmt}i,j,s,t). These regional  heterogeneities reflect differences in local sea ice dynamics, and circulation patterns that modulate how atmospheric forcing translates to surface
  buoyancy changes.


To understand the mechanisms underlying these SAM-WMT relationships, we now turn to examine the spatial patterns of surface heat flux and freshwater flux anomalies between high and low SAM composite years.

Heat flux composites reveal distinct responses between interglacial and glacial periods (Fig.~\ref{sam_heat}). The figure shows them as the surface density flux they produce, on the same scale as the freshwater composites in Fig.~\ref{sam_fwf}. During interglacials, positive SAM phases increase the ocean heat loss by more than 15~W~m$^{-2}$ over parts of the open ocean between 50 and 70°S, north of the winter ice edge (dashed lines in Fig.~\ref{sam_heat}a-c). These areas lie mainly in the Indian Ocean sector between 60 and 140°E, north of the Cooperation, Davis and Mawson Seas, and in the Pacific sector between 180° and 120°W, north of the Ross and Amundsen Seas, with smaller patches north of the Weddell Sea. Decomposition into components shows that turbulent fluxes (sensible and latent heat) account for most of this anomaly (Fig.~\ref{sam_heat}k-m), while radiative contributions (longwave and shortwave) remain small (Fig.~\ref{sam_heat}f-h). In contrast, glacial periods (LGM and MIS3) show only weak SAM-driven heat flux changes south of the winter ice edge, which lies near 55°S on average, because the extensive sea ice cover limits direct atmosphere-ocean exchange (Fig.~\ref{sam_heat}d,e).


Freshwater flux composites indicate that SAM influences dense water formation primarily through sea ice processes, and that net precipitation and river runoff contribute less (Fig.~\ref{sam_fwf}f-j against k-o). Total freshwater flux anomalies exhibit general dipole patterns with positive SAM driving freshwater loss (ocean salinification) in coastal regions accompanied by freshwater gain offshore (Fig.~\ref{sam_fwf}a-e). This pattern is controlled by sea ice processes, and the chain that produces it runs as follows. Strengthened and poleward-shifted westerlies associated with positive SAM increase the offshore component of the wind stress at the coast, which drives sea ice away from the continental margin. The divergence opens and widens coastal polynyas, exposing water already at the freezing point directly to a cold atmosphere. New ice therefore forms rapidly in those openings, and because ice rejects most of its salt as it grows, the salt is delivered to the surface ocean precisely where the water is coldest and densest. The ice exported offshore carries the compensating freshwater north, where it tends to melt and freshen the surface. This results in the observed dipole pattern, with salinification at the coast and freshening offshore. This is the same mechanism identified in observations of Antarctic Bottom Water production in the Cape Darnley polynya \cite{Ohshima2013} and in the sea ice contribution to Southern Ocean water mass transformation more generally \cite{Abernathey2016}.

These composites indicate that the mode acts on dense water formation through two routes. To separate them we computed, for the same high and low phases, the zonally integrated northward Ekman transport at 60°S and the coastal brine input integrated south of 65°S (Supplementary Fig. \ref{sam_pushpull}a,b). The positive phase strengthens the Ekman transport by 8.0--12.8~Sv in every climate state, which is the wind-driven upwelling that supplies water to the surface. It simultaneously increases the coastal brine input, by 51, 12 and 42~mSv in PI, the mid-Holocene and the Last Interglacial and by 24 and 31~mSv in LGM and MIS3, which is the salt pump that converts that water into denser classes. Both responses are positive in all five states, and so is the response of the peak transformation rate (Supplementary Fig. \ref{sam_pushpull}c).

The spatial patterns are shown in Supplementary Fig. \ref{sam_windice}. Under the positive phase the westerly wind stress strengthens over the whole circumpolar belt in every state (Supplementary Fig. \ref{sam_windice}a,c,e,g,i). The sea ice concentration changes mainly near the winter ice edge (Supplementary Fig. \ref{sam_windice}b,d,f,h,j). Gains and losses alternate between sectors along the edge, and the net winter ice area grows in all five states, by 0.07 to 0.57 million km$^2$. Near the coast the concentration hardly changes, because the winter ice there stays close to full cover. In the glacial states the ice edge response lies further north, following the expanded ice edge, and it is about twice as large as in the interglacial states.

Together, these results demonstrate that the Southern Annular Mode systematically modulates transformation rates in all five simulated climate states through two distinct contributors: thermal forcing, which dominates during interglacials via turbulent heat fluxes, and haline forcing, which dominates during glacials via sea ice processes.

Because Figs.~\ref{sam_heat} and \ref{sam_fwf} share one unit and one colour scale, the two forcings can be compared directly. Where sea ice is present, the freshwater term is the larger of the two over 90--94\% of the area, and on average it is 6 to 10 times larger than the heat term. North of the ice edge the heat term is the larger one over most of the area. This is the same split between ice-covered and open water that governs the mean state. These comparisons use local magnitudes. The freshwater term makes the surface denser near the coast and lighter further offshore, so its two parts partly offset each other in an area total.


%Climate state dependency is pronounced. Interglacial periods (PI/MH/LIG) exhibit spatially coherent positive heat flux anomalies throughout the ice-free Southern Ocean, with maximum values in the Atlantic and Indian sectors where open water permits direct atmosphere-ocean coupling. In contrast, glacial periods (LGM/MIS3) show reduced SAM-driven heat flux variability due to expanded sea ice insulation that dampens atmosphere-ocean exchange. The heat flux anomaly magnitude decreases by approximately 40\% during glacials compared to interglacials, demonstrating that sea ice extent fundamentally modulates SAM's capacity to influence surface buoyancy forcing.

%Regional patterns reflect local oceanographic conditions. The Ross Sea exhibits the strongest SAM sensitivity during interglacials, with heat flux anomalies exceeding 20~W/m$^2$ during high-SAM phases---consistent with this region's exposure to westerly wind variability and reduced sea ice buffering during warm periods. The Weddell Sea shows more muted response due to persistent open-water areas that maintain atmosphere-ocean coupling even during low-SAM conditions. This regional heterogeneity implies that SAM-driven WMT changes concentrate in specific sectors rather than occurring uniformly around Antarctica.


%High-SAM drives complex dipole patterns with coastal freshwater loss (salinification from enhanced brine rejection) and offshore gain (from ice melt). The sea ice component (row 2) dominates other sources (row 3) by factors of 3--5, confirming that SAM influences WMT primarily through haline forcing via sea ice redistribution. Glacial periods exhibit amplified responses ($>$1.0~mm/day) due to expanded baseline sea ice coverage providing larger reservoir for SAM-driven redistribution.


%The sea ice freshwater component dominates other sources (net precipitation and river runoff) in all climate states, confirming that SAM-driven WMT variability operates primarily through haline forcing via sea ice redistribution. Glacial periods exhibit amplified sea ice responses: LGM and MIS3 show freshwater flux anomalies exceeding 1.0~mm/day in coastal regions during high-SAM phases, compared to 0.5--0.7~mm/day in interglacials. This glacial amplification reflects expanded baseline sea ice coverage providing larger reservoir for SAM-driven redistribution.

%Seasonal dependence is critical. Winter composites (JJA) show strongest signals when sea ice formation peaks, whereas summer patterns (not shown) exhibit reversed but weaker anomalies associated with melt season dynamics. The asymmetry between winter intensification and summer weakening creates net annual haline forcing that persists beyond individual SAM events, potentially influencing dense water properties on multi-year timescales.


\subsection{Ventilation age simulations}


%\textbf{Figure 12: Vertical profiles of ideal age by ocean basin.}
%\begin{figure}[h]
%\centering
%\includegraphics[width=0.8\textwidth]{figures/age_vertical.pdf}
%\end{figure}
%Basin-averaged vertical age profiles (years vs depth) for Atlantic, Indian, and Pacific oceans across five experiments. Interglacials show modest depth-increasing ages: Atlantic/Indian 500--750~yr at 4000~m, Pacific 800--1200~yr. Glacials exhibit systematically older ages with strong depth dependence: Atlantic reaches 1500--2000~yr, Pacific exceeds 2500~yr below 3000~m (tripling relative to PI). Indian Ocean shows intermediate behavior. These profiles quantify basin-specific ventilation changes and demonstrate that Pacific deep waters experience most dramatic aging under glacial conditions due to remoteness from Southern Ocean formation regions.

%\textbf{Figure 13: Vertical profiles of ideal age anomalies by ocean basin.}
%\begin{figure}[h]
%\centering
%\includegraphics[width=0.8\textwidth]{figures/age_vertical_anm.pdf}
%\end{figure}
%Basin-averaged age anomaly profiles (years vs depth, relative to PI). MH exhibits minimal anomalies ($<\pm$100~yr). LIG shows slight rejuvenation in Atlantic ($-$75 to $-$150~yr throughout water column) from enhanced dense water formation, while Pacific shows modest increases (+100~yr deep). Glacials display depth-increasing anomalies: LGM/MIS3 reach +1500~yr in deep Pacific, +800~yr in Atlantic, +500~yr in Indian. The systematic depth-dependence reflects reduced ventilation at depth where sluggish abyssal circulation magnifies the impact of reduced formation rates. These anomaly profiles isolate pure climate-driven ventilation changes independent of baseline circulation structure.

\subsection{Overturning circulation and deep ocean ventilation}

The surface transformation changes described above should be reflected in the overturning circulation itself. We therefore examine the global overturning streamfunction computed from the simulated vertical velocity (Fig. \ref{moc}). In the interglacial states the abyssal cell, which represents the northward spreading of southern-sourced bottom water, reaches $-11.7$~Sv in PI, $-11.3$~Sv in the mid-Holocene and $-11.4$~Sv in the Last Interglacial. In the glacial states it weakens to $-2.3$~Sv in LGM and $-1.7$~Sv in MIS3, a reduction of roughly 80 to 86\%. The upper cell does not weaken and strengthens moderately, from 18.5~Sv in PI to 20.1~Sv in LGM and 21.5~Sv in MIS3. The glacial reorganisation in these simulations is therefore concentrated in the abyssal cell and expresses a contraction of the southern-sourced cell. The overturning as a whole remains active.

Having examined surface water mass transformation and the overturning response, we next assess deep ocean ventilation through ideal age distributions.


Ideal age distributions at 4000 m depth reveal fundamental restructuring of abyssal ventilation across different climate states (Fig.\ref{age}). Interglacial periods are characterized by relatively young deep waters. PI, the mid-Holocene and the Last Interglacial all present ages below 500 years in the Atlantic and Indian sectors, with the oldest waters (more than 500 years) confined to the Pacific.  The mid-Holocene exhibits no obvious age changes relative to PI. Last Interglacial anomalies reveal younger-than-PI ages (up to $-$150 years) in the Belingshausen Sea sector. This can be attributed to reduced sea ice coverage and enhanced thermal forcing. The Atlantic and Indian Ocean sectors show relatively modest age increases during the Last Interglacial (less than 70 years). Glacial periods exhibit different patterns. In LGM and MIS3, water ages exceed 750 year throughout the Atlantic and Indian basins, with Pacific ages reaching 1500 year. This represents a significant global increase in deep ocean isolation or a slowdown in overturning in glacial times, consistent with radiocarbon proxy reconstructions \cite{Skinner2017}.

We now examine ideal age throughout the global ocean (Fig. \ref{age_global}). Averaged below 2000 m, the age is strongly basin dependent. The Atlantic increases only modestly, from 377 years in PI to 439 years in LGM and 508 years in MIS3, because it remains ventilated from the north. The Southern Ocean more than doubles, from 381 to 789 and 844 years, and the Pacific, already the oldest basin in the interglacials at 847 years, reaches 1291 and 1314 years. The zonal-mean sections (Fig. \ref{age_global}f-j) show that the glacial age increase is concentrated below roughly 2000 m and is largest in the deep Pacific, while the upper ocean changes comparatively little. The contrast between a mildly ageing Atlantic and a strongly ageing Southern Ocean and Pacific is consistent with the abyssal cell weakening while the upper cell is maintained.

We caution that ideal age in these simulations should be read as a lower bound. Each experiment was integrated for 1000 years, which is shorter than the oldest simulated ages, so the glacial age fields are not fully equilibrated and the true equilibrium contrast between climate states would probably be larger. Achieving equilibrated deep-ocean age tracers requires multi-millennial integrations \cite{Millet2025}.


%Vertical age profiles quantify basin-specific ventilation changes (Fig.~12). In the Atlantic, glacial ages increase systematically with depth, reaching 1500--2000~yr at 4000~m compared to 500--750~yr in interglacials. The Indian Ocean shows intermediate behavior. Most striking is the Pacific, where glacial ages exceed 2500~yr below 3000~m---a tripling relative to PI. Age anomaly profiles (Fig.~13) isolate climate-driven changes: LIG shows slight rejuvenation throughout the water column in the Atlantic ($-$75 to $-$150~yr), while LGM/MIS3 exhibit depth-increasing anomalies reaching $+$1500~yr in the deep Pacific. These patterns reflect enhanced stratification and altered circulation under glacial boundary conditions.


% ============================================================================
% Methods
% ============================================================================


\section{Discussion and conclusions}

In this study, we analyse five simulated climate states spanning PI, the mid-Holocene, the Last Interglacial, LGM, and MIS3 to identify how surface buoyancy forcing mechanisms govern Southern Ocean dense water formation under varying boundary conditions. For this purpose, we apply the WMT framework \cite{Walin1982,Groeskamp2019,drake2025water}, which  provides a thermodynamic basis for relating surface buoyancy fluxes to dense water formation rates. Recent observational studies have successfully applied this framework to the modern Southern Ocean. For example, based on Argo floats and instrumented seals, Pellichero et al. \cite{Pellichero2018} show that seasonal sea ice growth and melt dominates WMT in the ice-covered region, with freshwater fluxes driving the meridional overturning that feeds bottom water formation, yielding an annual mean of approximately 5$\pm$5 Sv transformation into denser bottom water classes. Bailey et al. \cite{Bailey2023} further quantified regional WMT patterns, emphasizing the dominant role of haline forcing in the Weddell Sea. However, these observational WMT analyses are inherently limited to the modern climate state and sea ice regions. Our study extends the WMT framework to PI and paleoclimate conditions for the entire Southern Ocean, and decomposes total transformation into individual heat and salt components across glacial-interglacial time periods. For the PI state, our model reproduces comparable annual-mean transformation magnitudes to modern observations (Fig. \ref{s2}). The forcing balance, however, differs. Pellichero et al. \cite{Pellichero2018} and Bailey et al. \cite{Bailey2023} find haline dominance in the modern seasonal ice zone, whereas our PI integral is thermally dominated. Because this difference concerns the central claim of the present study, we examined it directly.

Two factors contribute, and they can be separated. The first is the integration domain. Pellichero et al. analyse only the region enclosed by the September sea ice edge, whereas we integrate over everything south of 60°S. In our PI simulation, 42.7\% of that area lies outside the September 15\% ice contour and is water they do not analyse. It is also the part of the domain where the ocean is exposed to the atmosphere all year and where thermal forcing is strongest, and taken alone it is 89--91\% thermally driven in the dense classes. Recomputing the transformation over a sea ice sector defined exactly as theirs, which reproduces their domain to within 2.5\% in area, moves the partition in the direction they report. The sea ice contribution at the transformation maximum roughly doubles, from 0.86 to 1.84~Sv (Fig. \ref{domain}).
% claude（11-27，0.1 格）：0.1 格下 Fig. 8 的峰值处海冰项是 2.25（60°S 以南）对 2.17 Sv（海冰区），不再翻倍，这句待你定。
 Resolved by density class, only one class of the PI transformation in this sector remains thermally dominated, $\sigma_2 = 36.75$--$37.00$~kg~m$^{-3}$, and every other class is haline dominated.

The second factor is the density range, and it accounts for most of the remaining difference. Converting our simulated surface properties inside the September ice zone to the neutral density coordinate used by Pellichero et al., our thermally dominated transformation maximum sits near $\gamma_n \approx 27.3$~kg~m$^{-3}$, whereas the haline-dominated lower cell they describe occupies $\gamma_n = 27.9$--$28.8$~kg~m$^{-3}$. Our surface transformation barely reaches the lighter edge of that range. In the densest classes our fluxes do populate ($\sigma_2 \ge 37.0$~kg~m$^{-3}$), the PI transformation is haline dominated, with 0.79~Sv from sea ice against 0.06~Sv from heat. The two analyses are therefore largely describing different water masses, ours weighted toward mode and intermediate densities and theirs toward bottom water densities, and where they overlap they agree in sign. This also indicates that the interglacial thermal pathway identified here is more relevant to the formation of intermediate and mode waters than to bottom water proper.

A residual discrepancy remains, and it reflects a genuine limitation of the model. We can exclude one candidate explanation: the simulated ice-covered sector is not warm biased in a way that would inflate the thermal term. Inside the September ice zone the area-weighted winter surface temperature is $-1.08$~°C, the median is $-1.71$~°C, 55\% of the area lies within 0.5~°C of the local freezing point, and the corresponding thermal expansion coefficient is $3.5\times10^{-5}$~K$^{-1}$. The near-freezing, weakly thermally sensitive regime invoked by Pellichero et al. is therefore reproduced. What differs is where dense water is made. In the real ocean, dense water forms through a sequence of ice shelf melt, dense shelf water formation in coastal polynyas, and downslope overflow \cite{Orsi1999,ToggweilerSamuels1995}. In the model, PI dense water forms by open-ocean convection in the interior of the Weddell gyre near 68°S. Because open-ocean convection exposes the water column to the atmosphere over a large area, it recruits an excessive thermal contribution and produces water that is not dense enough, which is precisely the offset in density class documented above.

To see where each density class is produced, we split the winter transformation south of 60°S by region (Supplementary Fig. \ref{formation_loc}). In the interglacial states the class of the transformation maximum forms mainly over ice-free deep water, which carries 68--71\% of it, while the continental shelf and slope, where the bottom is shallower than 1000~m, carry 16--25\%. The densest classes that each interglacial state produces carry 3.0--7.4~Sv and form mainly on the shelf and slope (61--88\%). The rest of them forms in the open water of the Weddell sector. Water as dense as the simulated PI bottom water is produced mostly offshore, with 13\% from the shelf. The model therefore does form dense water near the coast. In the interglacial states this coastal water is denser and much smaller in volume than the water made offshore. In the glacial states the class of the transformation maximum forms under sea ice offshore (87--90\%), and the densest classes form both on the shelf (53--59\%) and under the ice further offshore. Because the density of bottom water differs between the climate states, a single density threshold cannot isolate the source of bottom water across all five states. For this reason our analysis covers the transformation of all water south of 60°S. Surface transformation shows where water becomes dense, and it does not show where that water ends up.

Independent evaluations of this model provide context for these limitations, and we summarise the relevant points here. \citet{sidorenko2019evaluation} evaluate the same coupled configuration and report an abyssal overturning cell of about 10~Sv, comparable to the pre-industrial value obtained here. They also find that the Antarctic sea ice of this configuration is thin, hardly reaching 0.25~m in the central Weddell Sea in September, and that the Southern Ocean surface is too warm at this resolution, which they associate with the reduced sea ice area. Their overall assessment is that although biases in the representation of water mass properties and ventilation mechanisms are present, the model still produces a reasonable density distribution that maintains realistic transports. Our own comparisons with observation-based products point to the same bias. Compared with the ORAS5 reanalysis \cite{Zuo2019}, the winter mixed layer of the pre-industrial run is too deep in the interior of the Weddell gyre and too shallow along the Antarctic coast (Supplementary Fig. \ref{mld_obs}). Compared with the OSI SAF satellite sea ice record for 1979--2008 \cite{OSISAF2022}, the simulated September sea ice extent is about half of the observed value. Most of the deficit is in the Weddell sector, where the simulated ice edge lies about 10° of latitude too far south, and almost all of the ice melts in summer. A warm Southern Ocean with too little winter sea ice is the state in which open-ocean convection is favoured over shelf processes, so these biases are consistent with the convection behaviour we diagnose above.

AWI-ESM2 shares this bias with the current generation of climate models. In the CMIP6 ensemble, 28 of 35 models form deep and bottom water by open-ocean deep convection that is too deep, too frequent, or too extensive, and no model shows convincing shelf export in the Weddell Sea \cite{Heuze2021}; the same behaviour was documented in CMIP5 \cite{Heuze2013}. The present configuration also has no ice shelf cavities, so basal melt beneath floating ice and the associated production of near-freezing shelf water are absent. We state these limitations plainly because they bound what can be claimed. Our results should be read as a statement about how the partitioning of surface buoyancy forcing between heat and freshwater responds to changing sea ice cover, which is a robust property of the simulations. They do not give a quantitative reconstruction of past bottom water formation rates.

The shift remains when the analysis is restricted to the sea ice zone. There the glacial states are 15\% and 18\% thermally driven and the interglacial states 53--78\% (Fig. \ref{domain}), so the contrast does not come from the open water convection that the interglacial periods include. The shift also follows the sea ice cover. Where ice covers the ocean, it shields the surface from heat loss to the atmosphere, while brine is still released where new ice grows. The surface buoyancy forcing would become more brine-driven as the ice concentration increases.

Our conclusions about the glacial state are consistent with recent work using the same ocean model in a forced configuration. Chen et al. \cite{Chen2025} attribute the large glacial volume of Atlantic bottom water to a substantial increase in sea ice export toward lower latitudes, which supplies dense shelf water year round, and to weaker mixing between northern- and southern-sourced water, which allows that water to retain the properties of its origin. The agreement is worth stating explicitly, and the two studies complement each other. Chen et al. infer the role of sea ice from water mass volumes, the overturning streamfunction, and mixed-layer and sea ice fields, in ocean-only simulations of two states with a prescribed atmosphere. The transformation framework applied here measures the surface buoyancy forcing directly and attributes it to individual flux components in sverdrups, so the contribution of sea ice is quantified. Running the analysis across five coupled climate states, in which sea ice and the atmospheric fluxes evolve together, is also what makes the change of regime visible, since a shift between two forcing regimes cannot be identified from a single glacial state. The present analysis therefore tests and quantifies the mechanism that Chen et al. proposed.

Our analysis addresses the surface forcing of all water masses south of 60°S. These results still have important implications for understanding how deep water formation mechanisms may respond to climate change (Fig. \ref{schematic}). Our finding that glacial sea ice acts as an amplifier of haline forcing aligns with the sea ice pump hypothesis \cite{Abernathey2016}. In this regime, the decoupling of thermal forcing implies that ventilation rates become sensitive mainly to sea ice production. Quantitatively, the interglacial Southern Ocean is characterized by thermally dominated WMT, with transformation maxima of 70-81~Sv at density classes of $\sigma_2 = 36.4$-$36.8$ kg m$^{-3}$, where turbulent heat exchange (sensible and latent heat fluxes) accounts for the majority of thermal forcing. Glacial conditions fundamentally restructure this forcing balance. The winter transformation maximum weakens to below 60~Sv and moves to $\sigma_2 = 37.1$--$37.6$~kg~m$^{-3}$, and the thermal share of the transformation falls to 13--18\% while the sea ice term becomes dominant. Extensive ice cover insulates the ocean from atmospheric heat exchange, and simultaneously intensifies haline forcing through enhanced brine rejection concentrated in coastal regions. This shift of forcing mechanisms can change pathways and distribution of dense water formation. Moreover, our paleoclimate analysis suggests that projected sea ice decline under global warming may lead to increased thermal dominance in the Southern Ocean.  However, concurrent Antarctic ice sheet mass loss could substantially modify this response through enhanced freshwater input \cite{Li2023Nature,Silvano2018}. The combined effects of reduced sea ice and increased meltwater highlight the need for comprehensive  assessments that account for both factors when projecting future changes in Southern Ocean ventilation and carbon uptake.

Many PMIP models underestimate glacial Southern Ocean sea ice and its seasonality \cite{Lhardy2021}. In our simulations the LGM sea ice extent is 29.8 million km$^2$ in winter and 11.1 million km$^2$ in summer. Both lie within the proxy-based estimates of about $33\pm5$ and $10\pm3$ million km$^2$ \cite{Lhardy2021}. We have not tested why the model does well here. The relatively fine ocean mesh near Antarctica may help, and the simulated glacial climate is cold enough to keep the surface water around Antarctica at its freezing point.

A related limitation concerns the simulated glacial winds. The westerlies in our glacial experiments shift poleward and strengthen, whereas proxy reconstructions infer an equatorward shift of about 4.8° and a weakening of about 25\% at the LGM relative to the mid-Holocene \cite{Gray2023}. Gray et al. \cite{Gray2023} found that the inferred shift exceeds that of any PMIP3/4 model. In the four PMIP4 LGM simulations \cite{Kageyama2021}, the winter wind maximum moves by less than 1.2° between LGM and PI, so current models in general do not reproduce the reconstructed change. Our model departs further in the opposite direction, with a poleward shift of 3.7°. Moreover, Gray et al. \cite{Gray2023} found no correlation across the model ensemble between Antarctic sea ice extent and either the latitude of the westerlies or the position of the surface temperature front. A bias in the position of the westerlies therefore need not carry over into the sea ice field on which our main result rests. Thus the SAM results describe how the model responds to a wind anomaly added to a background state whose mean winds are biased.

The Southern Annular Mode emerges as a persistent modulator of transformation across all climate states. In its positive phase it enhances turbulent heat loss over open water and increases brine rejection near the coast, with composite anomalies of 3--8 Sv between high and low phases. The two effects act together. A positive anomaly strengthens the Ekman divergence that brings deep water to the surface, and in the same winters it strengthens the coastal salt pump that converts that water into denser classes, so the wind acts on both the supply of water to the surface and its conversion once there. This circumpolar picture does not hold uniformly in every sector. The Ross Sea in particular departs from it, showing weaker and at some densities opposing responses, and in the glacial states the circumpolar signal is carried largely by the Weddell sector while the Ross and Ad\'elie sectors contribute little. The coupling holds in all five states, although their background conditions differ strongly. The size of the Ekman and brine responses does not separate the interglacial from the glacial states, and the transformation anomaly is of similar size in every state. What changes is the local balance of the anomalous forcing. Over open water the heat term dominates, and under sea ice the freshwater term dominates, so the expanded glacial ice cover shifts the response towards brine release in the same way as it shifts the mean state. This finding agrees with recent observational evidence linking Weddell Sea bottom water reduction to wind-driven sea ice decline \cite{Zhou2023}, and the recovery of Ross Sea convection during 2015--2019 from positive SAM combined with El Ni\~no forcing \cite{Silvano2020}. Whether this coupling could be used for proxy reconstruction is unclear. Recovering an interannual-to-decadal signal from marine sediments would require accumulation rates and chronological precision that are rarely available around the Antarctic margin, where sedimentation is slow and bioturbation mixes the record over intervals far longer than the variability of interest. We therefore present the persistence of the coupling as a property of the simulated system. We do not expect it to be recoverable from marine sediments.

Finally, we evaluate the simulated deep ocean ventilation ages. Benthic-planktonic $^{14}$C age offsets indicate that the radiocarbon age of dissolved inorganic carbon in the glacial deep ocean increased globally, with the strongest signals in the Pacific and Southern Ocean \cite{Skinner2010,Skinner2017}. Our  simulated LGM ideal age increase of about 1000 years at 4000 m depth is of similar magnitude to radiocarbon-derived estimates of  689$\pm$53~$^{14}$C-yr \cite{Skinner2017}, though our model somewhat overestimates the age increase. However, direct comparison requires careful interpretation because benthic-planktonic $^{14}$C ages include both circulation age and preformed age arising from incomplete atmosphere-ocean equilibration \cite{Koeve2015}. Our glacial simulations present substantially increased ventilation ages, with Pacific ages reaching 1500 years compared to less than 500 years during interglacials, consistent with radiocarbon proxy reconstructions \cite{Skinner2017}. Two caveats apply to this comparison. The ideal age fields are not fully equilibrated after 1000 years of integration, so the simulated glacial ages, and the glacial-interglacial contrast, are best read as lower bounds; equilibrating deep-ocean age tracers requires multi-millennial simulations \cite{Millet2025}. Ideal age also responds to mixing as well as to advection, so it measures the combined effect of a weaker abyssal cell and altered interior mixing. The overturning streamfunction shown in Fig. \ref{moc} provides the more direct measure, and the two diagnostics agree in indicating a substantially weaker glacial abyssal cell.

While this study is based on a single model, the mechanisms identified here provide a framework for systematic inter-model comparison. Applying this analysis across the CMIP/PMIP ensemble would help address the current large inter-model spread in glacial-interglacial WMT and better constrain its implications for deep water formation.

These paleoclimate results carry direct implications for the future Southern Ocean. The thermal-to-haline transition identified here reflects a general state-dependence of surface buoyancy forcing. The same competing mechanisms are projected to operate, in the opposite direction, under anthropogenic warming. Sea-ice decline reduces the insulation that currently suppresses open-ocean heat loss, favouring a more thermally driven regime and, at the extreme, the cessation of open-ocean deep convection \cite{deLavergne2014}. Acting against this, accelerating Antarctic ice-sheet mass loss freshens the surface and weakens the haline forcing that sustains dense water formation in coastal polynyas \cite{Silvano2018,Jeong2023}, a process already implicated in the observed contraction, warming, and reduced ventilation of the abyssal Southern Ocean \cite{Li2023Nature,Gunn2023}. The net outcome depends on which of these two pathways dominates, and our analysis shows that this balance is itself climate-state dependent. This implies that projections of Southern Ocean ventilation and its associated carbon uptake will remain uncertain until the thermal and haline components of surface buoyancy forcing are evaluated separately, alongside bulk overturning metrics.

\section{Methods}

\subsection{Model description and experimental design}

We employ the AWI-ESM2, a state-of-the-art coupled Earth system model developed at the Alfred Wegener Institute \cite{sidorenko2019evaluation}. The model comprises the ECHAM6 atmospheric component \cite{stevens2013atmospheric} coupled to the JSBACH land-surface model, which represents dynamic vegetation through multiple plant functional types \cite{brovkin2009global,reick2021jsbach}. The ocean-sea ice component is FESOM2, which features a multi-resolution dynamical core based on a finite-volume formulation \cite{danilov2017finite}. The atmospheric component operates at T63L47 spectral resolution (approx. 1.875° horizontal resolution with 47 vertical levels). A key feature of AWI-ESM2 is the unstructured mesh of the ocean component, allowing for spatially variable resolution (Fig. \ref{reso}). In our configuration, the grid resolution ranges from $\sim$100 km in the open ocean to 25 km in polar regions and along coastlines, with 46 distinct vertical layers. At this resolution the larger coastal polynyas are represented as features, but the boundary layer processes that set polynya dynamics are not resolved. Sea ice growth and melt are computed thermodynamically at each surface node, and the associated brine release enters the ocean as a salt flux at that node, so the brine signal appears wherever the model forms ice. Its intensity depends on the simulated wind-driven ice divergence, since the mesh does not resolve the small-scale processes inside the polynyas. Resolving them would need a considerably finer mesh near the continent, of the order of a few kilometres. The configuration does not include ice shelf cavities, so the ocean beneath floating ice shelves, and the basal melt it supplies, are absent from the simulations. Regions covered by floating ice shelves are therefore outside the model domain and appear blank in the map figures.
  The ocean component employs the K-Profile Parameterization (KPP) scheme
  \cite{large1994oceanic} for vertical mixing, which detects the ocean boundary
  layer depth using a bulk Richardson number criterion ($Ri_{cr} = 0.3$) and
  includes a non-local transport term to represent convective fluxes within the
  unstable boundary layer. Static gravitational instabilities throughout the water
  column are treated by enhancing the vertical diffusivity to 0.1~m$^2$~s$^{-1}$
  wherever the squared Brunt-V\"{a}is\"{a}l\"{a} frequency becomes negative.
  In addition, a Monin-Obukhov length scale mixing scheme following
  \cite{timmermann2004parameterization} is applied at high latitudes to represent
  surface buoyancy flux-driven convection, with a mixing coefficient of
  0.01~m$^2$~s$^{-1}$. Mesoscale eddy-induced tracer transport is represented through
  Gent-McWilliams thickness diffusion \cite{gent1990isopycnal,ferrari2008parameterization}
  coupled with Redi isopycnal diffusion \cite{redi1982oceanic}.  The GM diffusivity
  is allowed to vary between 2 and 2000~m$^2$~s$^{-1}$, scaled according to the
  local model resolution relative to a reference scale of 100~km.

We perform five equilibrium simulations spanning distinct climate states including three interglacial periods, i.e., pre-industrial (PI), mid-Holocene (MH), and Last Interglacial (LIG), and two glacial periods, i.e., the Last Glacial Maximum (LGM) and Marine Isotope Stage 3 (MIS3). The boundary conditions are configured following the criteria of PMIP4 \cite{otto2017pmip4} and are summarised in Table \ref{tab:bc}. Orbital parameters are calculated following \cite{berger1977long}, and the greenhouse gas concentrations are taken from multi-archive reconstructions from ice core records and recent measurements of firn air and atmospheric samples \cite{fluckiger2002high,monnin2004evidence,schilt2010glacial,buiron2011taldice,schneider2013reconstruction,kohler2017156}.  For LGM and MIS3, we fix the boundary conditions at 21 ka and 38 ka respectively, with the topography and ice-sheet properties deriving from the GLAC1D reconstruction \cite{tarasov2003greenland,tarasov2012data,briggs2014data}. The GLAC-1D reconstruction enters the model through several boundary fields at once, and it is clearest to list them by the model component that uses them. Specifically, the atmospheric model component receives the ice sheet extent and surface elevation as orography, together with the land-sea mask implied by the lowered sea level. The land surface model component receives the same topography, from which river routing and vegetation distribution are derived. The ocean model component receives bathymetry and a modified coastline. Because the ocean component uses an unstructured mesh, a changed coastline cannot be imposed on the existing grid and a dedicated mesh must be generated for each glacial period. We use meshes constructed for 21 ka and 38 ka respectively, while PI, the mid-Holocene and the Last Interglacial share the modern-geometry mesh. Global mean ocean salinity is raised in the glacial experiments to account for the water stored in the ice sheets. No prescribed ice sheet meltwater flux is applied in any of the equilibrium simulations, and the model does not include ice shelf cavities, so basal melt beneath floating ice shelves is not represented. Precipitation minus evaporation over the Antarctic ice sheet is routed to the coast by the hydrological discharge model and enters the ocean as liquid runoff, about 0.085~Sv south of 60°S in PI and 0.07~Sv in the glacial states. This runoff stands in for the freshwater that the ice sheet would otherwise lose by basal melt and iceberg calving.

The PI control run was integrated for 1,500 years, initialized with AMIP atmospheric conditions \cite{roeckner2004atmospheric} for the atmosphere and World Ocean Atlas climatology \cite{levitus2010world} for the ocean. The mid-Holocene and Last Interglacial simulations branched from the equilibrated PI state. The LGM and MIS3 experiments were initialized from a previous glacial state \cite{werner2016glacial}. All paleo-simulations were integrated for 1,000 model years. This is the same ensemble, with the same spin-up and the same boundary conditions, that is described in the experimental design of \citet{Shi2025GRL}. We analyze the final 100 years of each run, representing a quasi-equilibrium state. Over these 100 years the trend of global mean sea surface temperature is within $\pm0.03$~K~century$^{-1}$ in PI, the mid-Holocene, the Last Interglacial and LGM and $+0.06$~K~century$^{-1}$ in MIS3. The trend of the Atlantic meridional overturning circulation is not significant, except for $+1.3$~Sv~century$^{-1}$ in PI ($p = 0.04$). These are the two quantities by which the PMIP4 protocol judges equilibrium \cite{Kageyama2017}.

These simulations have been documented and evaluated against proxy data in previous studies with this model, and the five-member ensemble analysed here is the same one used in those studies. Because the balance of agreement and disagreement in those evaluations bears on how far the present results can be trusted, we summarise it here. For the pre-industrial state, \citet{Shi2025GRL} find precipitation in reasonable agreement with the GPCP observations, with local wet biases at the position of the Intertropical Convergence Zone in both solstitial seasons, a bias common to other coupled models. For the interglacials, \citet{Shi2022JCLI} compare both generations of AWI-ESM against the pollen compilation of \citet{Bartlein2011} and the multi-archive compilation of \citet{TurneyJones2010}. The model reproduces the mid-Holocene warming over western Europe and the north-south dipole of the temperature anomaly over eastern North America, and it captures the sign and magnitude of the wetting over the southern Sahara and South Asia. It disagrees with the proxies over North Africa, where the simulated mid-Holocene is cooler than the pre-industrial while the reconstructions show warming, and it underestimates the mid-Holocene wetting over the northern Sahara. For the Last Interglacial the simulated Northern Hemisphere warming and the tropical and North African cooling both match the records, but the Antarctic warming that the proxies show is absent, which \citet{Shi2022JCLI} attribute to the PMIP4 protocol holding the Last Interglacial ice sheet orography fixed at its pre-industrial configuration. For the glacial states, \citet{Shi2023CP} document the Last Glacial Maximum configuration against the Bartlein pollen compilation and MARGO sea surface temperatures, and \citet{Shi2025GRL} find that model and proxies agree on the pronounced drying over most Northern Hemisphere continents at the Last Glacial Maximum, while for Marine Isotope Stage 3 no comprehensive precipitation compilation exists and the comparison rests on individual sites. Across all four paleoclimate states the simulated anomalies are consistent with the PMIP4 ensemble. We note that the present analysis depends on Southern Ocean hydrography, sea ice and overturning, and none of these evaluations examines these fields against observations. They therefore constrain only the large-scale climate response of the model.


\subsection{Surface buoyancy flux decomposition}

To quantify the individual contributions of different physical processes to surface buoyancy forcing, we employ the xbudget diagnostic tool (\url{https://github.com/hdrake/xbudget}) \cite{drake2025water}. This framework enables systematic decomposition of tracer budgets into individual forcing terms while maintaining numerical consistency with the underlying model's finite-volume discretization.

The total surface density tendency, $\mathcal{B}$ (kg m$^{-2}$ s$^{-1}$; divided by 1 m it gives the density change of the top metre of the water column, and in the figures it is shown per month of 30 days), is decomposed into thermal ($\mathcal{B}_{therm}$) and haline ($\mathcal{B}_{haline}$) components:
\begin{equation}
\mathcal{B} = \mathcal{B}_{therm} + \mathcal{B}_{haline} = -\frac{\alpha}{C_p} Q_{net} - \beta S F_{fw},
\end{equation}
where $\alpha$ and $\beta$ are the thermal expansion and haline contraction coefficients, respectively, and $C_p$ is the specific heat capacity of seawater. $Q_{net}$ represents the net surface heat flux (positive into the ocean). Here we further decompose the net heat flux into its shortwave radiative, longwave radiative, latent heat, and sensible heat components. In the figures the radiative component denotes the sum of the shortwave and longwave fluxes, and the turbulent component denotes the sum of the latent and sensible heat fluxes. Similarly, the freshwater fluxes $F_{fw}$, taken as mass fluxes in kg m$^{-2}$ s$^{-1}$ and positive into the ocean, are partitioned into precipitation, evaporation, river runoff, and thermodynamic sea ice processes (brine rejection and freshwater release upon melting). This systematic decomposition allows us to isolate the specific role of sea ice thermodynamics versus atmospheric radiative and turbulent fluxes in driving water mass transformation. Under this convention, a heat loss ($Q_{\text{net}} < 0$) or a freshwater loss ($F_{fw} < 0$) gives a positive $\mathcal{B}$, that is, densification.

We refer throughout to the thermal density tendency, meaning the surface density change produced by the net heat flux acting through the thermal expansion coefficient, that is the $-\left(\alpha/C_p\right)Q_{net}$ term, and to the sea-ice-driven density tendency, meaning the part of the haline term $-\beta S F_{fw}$ that is carried by the thermodynamic growth and melt of sea ice. A positive value denotes densification of the surface water in both cases. Brine released during ice growth increases surface salinity and therefore density, whereas melting releases freshwater and reduces it.

After xbudget decomposes surface density tendency into individual heat and salt components, these decomposed budgets are passed directly to xwmt (see the next subsection) for WMT calculation.


\subsection{Water mass transformation framework}

We quantify the production of dense water using the water mass transformation (WMT) framework \cite{Walin1982,Speer1992}. This method relates surface density tendency to the cross-isopycnal volume flux, providing a mechanistic link between air-sea-ice forcing and the overturning circulation.

The water mass transformation rate, $F(\rho)$, represents the volume flux across a specific isopycnal surface $\rho$ driven by surface forcing. It is calculated by integrating the buoyancy flux over the outcrop area where the surface density is less than or equal to $\rho$:
\begin{equation}
F(\rho) = \iint_{\rho_{surf} \le \rho} \frac{\mathcal{B}}{\Delta \rho} \, dA,
\end{equation}
where $A$ denotes the surface area. In the discrete formulation used for analysis, we bin the surface fluxes into density classes with a bin width ($\Delta \rho$) of 0.1 kg m$^{-3}$.

Finally, the water mass formation rate, $M(\rho)$, is defined as the convergence of the transformation rate in density space:
\begin{equation}
M(\rho) = - \frac{\partial F(\rho)}{\partial \rho}.
\end{equation}
A positive $M(\rho)$ indicates that surface fluxes act to accumulate water in a specific density class, effectively transferring water from the surface mixed layer into the interior ocean (subduction/ventilation). By decomposing $\mathcal{B}$, we can explicitly attribute changes in dense water formation to either thermal or haline forcing.

We calculate these diagnostics using the xwmt package \cite{xwmt2024}, a Python-based framework developed at NOAA-GFDL specifically designed for calculating WMT in a computationally efficient manner \cite{drake2025water}. For each experiment, WMT is computed over four regions south of 60°S: (1) the entire Southern Ocean, (2) the Ross Sea sector (180°W--60°W), (3) the Weddell Sea sector (60°W--79°E), and (4) the Adélie sector (79°E--180°E). These regions encompass the primary Antarctic bottom water (AABW) formation sites and allow regional comparison of formation mechanisms. Regional transformations are calculated by applying spatial masks to the surface forcing fields before integration.

To compare with observational estimates that are defined over the ice-covered sector, we additionally repeat the calculation over a seasonal sea ice zone, defined as the region where the climatological September sea ice concentration exceeds 15\%, following the domain definition of Pellichero et al. \cite{Pellichero2018}, and over the complementary part of the domain south of 60°S that lies outside that contour. These auxiliary calculations use the climatological monthly means and are otherwise identical to the main calculation. For the density-class comparison with observational studies that use neutral density, simulated surface properties were converted with the TEOS-10 routines.

\subsection{Overturning streamfunction}

The global meridional overturning streamfunction is computed from the simulated vertical velocity on the ocean model mesh. The vertical volume transport across each model level is integrated over the triangular mesh elements and accumulated from the north in 2° bins of geographic latitude. Annual means are averaged over the last 20 years of the 100-year analysis period of each simulation. The strength of the abyssal cell is defined as the minimum of the streamfunction below 2000~m and south of 30°S. The strength of the upper cell is defined as the maximum between 500 and 4000~m depth and between 20°N and 70°N.

\begin{table}[htbp]
\centering
\caption{Boundary conditions prescribed in the five equilibrium simulations. Orbital parameters for PI are those of 1850~CE. Greenhouse gas concentrations are volume mixing ratios.\label{tab:bc}}
\begin{tabular}{lccccccc}
\toprule
 & CO$_2$ & CH$_4$ & N$_2$O & Eccentricity & Obliquity & Perihelion & Ice sheet / \\
 & (ppm) & (ppb) & (ppb) & & (°) & (°) & bathymetry \\
\midrule
PI   & 284.3 & 808.2 & 273.0 & 1850~CE  & 1850~CE & 1850~CE & modern \\
MH   & 264.4 & 597.0 & 262.0 & 0.018682 & 24.105  & 180.87  & modern \\
LIG  & 275.0 & 685.0 & 255.0 & 0.039378 & 24.040  & 95.41   & modern \\
LGM  & 190.0 & 375.0 & 200.0 & 0.018994 & 22.949  & 294.42  & GLAC-1D 21~ka \\
MIS3 & 210.5 & 556.2 & 247.4 & 0.013676 & 23.2591 & 25.99   & GLAC-1D 38~ka \\
\botrule
\end{tabular}
\end{table}


% ============================================================================
% Discussion and Conclusions
% ============================================================================

% ============================================================================
% Backmatter
% ============================================================================
\backmatter

\section*{Acknowledgements}
%% TODO before submission: acknowledgements and funding text.
%% TODO before submission: verify the Sidorenko et al. (2019) statements in the
%%   Discussion (about 10 Sv abyssal cell, 0.25 m September Weddell ice thickness,
%%   too-warm Southern Ocean surface in the historical run, and the "reasonable
%%   density distribution" assessment) against the published PDF. The 1.43 K value
%%   (global SST RMSE, not a Southern Ocean bias) was removed on 2026-11-25. No copy of that paper is on the
%%   analysis machine and the project notes disagree about whether the numbers
%%   were ever checked against the full text, so they are unconfirmed here.
This research was supported by... We acknowledge helpful discussions with...

\section*{Author contributions}
X.S. and J.L. designed the study, performed the analysis, and wrote the manuscript. All authors discussed the results and commented on the manuscript.

\section*{Competing interests}
The authors declare no competing interests.

\section*{Data availability}
The model output required to reproduce every figure in this manuscript, comprising climatological monthly means for the five experiments, the 100-year monthly water mass transformation time series, the Southern Annular Mode composite fields, and the ocean mesh files, has been deposited in a public repository (DOI to be inserted upon acceptance). Source data for the line-graph figures are provided with this paper. The raw model output totals several tens of terabytes and is archived at the German Climate Computing Centre; it is available from the corresponding author on reasonable request.

\section*{Code availability}
The xbudget tool used for surface density tendency decomposition is available at \url{https://github.com/hdrake/xbudget}. The xwmt package used for water mass transformation calculations is available at \url{https://github.com/NOAA-GFDL/xwmt}. The analysis and plotting scripts specific to this study, covering the surface buoyancy flux decomposition, the transformation calculations including the sea ice sector comparison, the Southern Annular Mode compositing, and all figures, are archived together with the data deposit and are also available at \url{https://github.com/xshi-awi/aabw_5exps}.

\bibliographystyle{sn-nature}
\bibliography{ref}

\clearpage


\begin{figure}
\centering
\includegraphics[width=\textwidth]{figures/fig_main_climate_jja.pdf}
 \caption{Simulated June-September mean anomalies in the mid-Holocene, the Last Interglacial, LGM, and MIS3 relative to PI for (a-d) sea surface temperature   [$^{\circ}$C], (e-h) sea surface salinity  [psu],  (i-l) sea ice concentration, and (m-p) 10 m zonal wind  [m/s]. Vectors represent PI wind. Four columns represent the mid-Holocene, the Last Interglacial, LGM, and MIS3, respectively. \label{sst}}
\end{figure}


\begin{figure}
\centering
\includegraphics[width=\textwidth]{figures/plot_wmt_4regions_5exps_winter_from_100years.pdf}
\caption{JJA-mean water mass transformation rates [Sv] as a function of density $\sigma_2$ [kg m$^{-3}$], in density bins 0.1 kg m$^{-3}$ wide, for (a-e) Southern Ocean, (f-j) Ross Sea, (k-o) Weddell Sea and (p-t) Adélie Land. Five columns represent five climate experiments (PI, mid-Holocene, Last Interglacial, LGM, and MIS3).  Lines show the mean values; shading denotes ±1  standard deviation across years.    \label{wmt_jja}}
\end{figure}


\begin{figure}
\centering
\includegraphics[width=\textwidth]{figures/surface_heat_density_tendency_from_raw_jja.pdf}
\caption{Simulated JJA-mean surface density tendency caused by (a-e)  total heat flux, (f-j) radiative fluxes (shortwave plus long wave fluxes), and (k-o) turbulent fluxes (latent and sensible fluxes).  Column 1 shows PI absolute values, and columns 2-5 show anomalies in the mid-Holocene, the Last Interglacial, LGM and MIS3 relative to PI, respectively.  Units: kg m$^{-2}$ month$^{-1}$, a surface density flux, which equals the density change of a surface layer 1 m thick over one month. \label{density_heat}}
\end{figure}


\begin{figure}
\centering
\includegraphics[width=\textwidth]{figures/surface_freshwater_density_tendency_from_raw_jja.pdf}
\caption{Simulated JJA-mean surface density tendency caused by (a-e)  total freshwater flux, (f-j) sea ice-induced freshwater flux, and (k-o) net precipitation plus river runoff. Column 1 shows PI absolute values, and columns 2-5 show anomalies in the mid-Holocene, the Last Interglacial, LGM and MIS3 relative to PI, respectively. Units: kg m$^{-2}$ month$^{-1}$, a surface density flux, which equals the density change of a surface layer 1 m thick over one month. \label{density_fwf}}
\end{figure}




\begin{figure}
\centering
\includegraphics[width=\textwidth]{figures/figR2_moc_5exps.pdf}
\caption{Global meridional overturning streamfunction computed from the simulated vertical velocity and averaged over the last 20 years of the 100-year analysis period. (a-e) Absolute streamfunction for PI, the mid-Holocene, the Last Interglacial, LGM and MIS3. Positive values (red) denote clockwise circulation in the latitude-depth plane; the negative (blue) cell below about 2500 m represents the northward spreading of southern-sourced bottom water. The strength of that abyssal cell, defined as the streamfunction minimum below 2000 m and south of 30°S, is given in each panel. (f-i) Anomalies relative to PI for the mid-Holocene, the Last Interglacial, LGM and MIS3. Units: Sv. \label{moc}}
\end{figure}

\begin{figure}
\centering
\includegraphics[width=\textwidth]{figures/age_horizontal_3900m_combined.pdf}
\caption{(a-e) Simulated spatial distribution of Southern Ocean ventilation age at 4000 m depth for (a) PI, (b) mid-Holocene, (c) Last Interglacial, (d) LGM,  (e) MIS3.  (f-i) Simulated anomalies of ventilation age at 4000 m for (f) mid-Holocene minus PI, (g) Last Interglacial minus PI, (h) LGM minus PI, and (i) MIS3 minus PI. Units: year. \label{age}}
\end{figure}

\begin{figure}
\centering
\includegraphics[width=\textwidth]{figures/figR3_age_global.pdf}
\caption{Global ideal age. (a-d) Basin-mean vertical profiles for the Atlantic, Indian, Pacific and Southern Ocean, for the five climate states. (e) Basin-mean ideal age below 2000 m. (f-j) Zonal-mean ideal age sections for PI, the mid-Holocene, the Last Interglacial, LGM and MIS3. Basins are defined as Atlantic 70°W-20°E and 35°S-65°N, Indian 20°E-115°E and 35°S-25°N, Pacific 120°E-70°W and 35°S-65°N, and Southern Ocean south of 35°S. Units: years. \label{age_global}}
\end{figure}

\begin{figure}
\centering
\includegraphics[width=\textwidth]{figures/figR1_wmt_domain_sensitivity.pdf}
\caption{Sensitivity of the thermal and haline partition to the integration domain, shown as annual-mean transformation rates for the five climate states. (a-e) Integrated south of 60°S, the domain used throughout this study. (f-j) Integrated over the seasonal sea ice zone only, defined as the region where the climatological September sea ice concentration exceeds 15\%, which reproduces the domain of Pellichero et al. \cite{Pellichero2018} to within 2.5\% in area. The share of the total transformation carried by the thermal term at the transformation maximum, and the integrated area of each domain, are given below each panel. Restricting the integration to the ice-covered sector increases the haline share, as expected, but leaves the contrast between thermally dominated interglacials and haline dominated glacials intact. Units: Sv. \label{domain}}
\end{figure}




\begin{figure}
\centering
\centering\includegraphics[angle=270,width=\textwidth]{figures/frame.pdf}
\caption{Mechanistic regime shift between interglacial and glacial Southern Ocean dense water formation. Two-panel schematic illustrating the thermal-to-haline transition in surface buoyancy forcing across climate states. Left: under interglacial conditions (PI, mid-Holocene, Last Interglacial), limited sea ice cover allows strong turbulent heat loss over the open ocean to drive surface cooling and dense water formation at $\sigma_2 \approx 36.5$--$36.8$~kg~m$^{-3}$ (60--80~Sv). Right: under glacial conditions (LGM, MIS3), expanded sea ice insulates the ocean from atmospheric heat exchange while intensifying brine rejection in coastal polynyas, producing denser waters ($\sigma_2 > 37.0$~kg~m$^{-3}$) at lower rates (48--58~Sv). Ventilation ages at 4000~m depth correspondingly shift from $<500$~yr (interglacial) to $>750$~yr basin-wide (glacial), with Pacific values reaching $\sim$1500~yr. Positive SAM phases enhance dense water formation across all states, primarily through amplified heat loss in interglacials and amplified brine rejection in glacials. \label{schematic}}
\end{figure}

    
\clearpage

\begin{appendices}

\section{Supplementary figures}\label{secA1}

\begin{figure}
\centering
\includegraphics[width=\textwidth]{figures/figS_density_windstress_jja.pdf}
\caption{Simulated June-September mean anomalies in the mid-Holocene, the Last Interglacial, LGM, and MIS3 relative to PI for (a-d) sea surface density [kg/m$^3$], and (e-h) 10 m zonal wind stress [N/m$^2$]. Vectors represent PI wind stress. Four columns represent the mid-Holocene, the Last Interglacial, LGM, and MIS3, respectively. \label{wind}}
\end{figure}

\clearpage

\begin{figure}
\centering
\includegraphics[width=\textwidth]{figures/fig02_climate_mld_seaice_jja_mld1000.pdf}
\caption{Simulated June-September mean (a-e)  mixed layer depth [m] and (f-j) sea ice concentration for (a,f) PI, (b,g) mid-Holocene, (c,h) Last Interglacial, (d,i) LGM, and (e,j) MIS3. The mixed layer depth colour scale runs to 1000 m in (a-c) and to 400 m in (d,e). \label{mld}}
\end{figure}

\clearpage

\begin{figure}
      \centering\includegraphics[width=\textwidth]{figures/antarctic_resolution_wireframe.pdf}
        \caption{Spatial resolution of the Southern Ocean. \label{reso}}
    \end{figure}

\clearpage

\begin{figure}
      \centering\includegraphics[width=\textwidth]{figures/figR7_mld_glacial_zoom.pdf}
        \caption{Simulated JJA-mean mixed layer depth south of 60°S for (a) LGM and (b) MIS3, on a 0--300 m colour scale. Units: m. \label{s1}}
    \end{figure}

\clearpage

\begin{figure}
\centering
\includegraphics[width=\textwidth]{figures/figS_wmt_common_axis_winter.pdf}
\caption{As Fig. \ref{wmt_jja}, but with a common vertical and horizontal axis for all panels, so that transformation magnitudes can be compared directly between climate states and regions. Rows show the Southern Ocean, Ross Sea, Weddell Sea and Ad\'{e}lie Land sectors; columns show PI, the mid-Holocene, the Last Interglacial, LGM and MIS3. Lines show 100-year means and shading denotes $\pm$1 standard deviation across years. Units: Sv. \label{s_common_axis}}
\end{figure}

\clearpage

\begin{figure}
\centering
\includegraphics[width=\textwidth]{figures/wmt_sam_composite_4regions_5exps.pdf}
\caption{Composite JJA-mean anomalies of water mass transformation rates [Sv] as a function of potential density $\sigma_2$ [kg m$^{-3}$] for high-SAM minus low-SAM years ($\pm$1.2$\sigma$ threshold), shown for (a-e) Southern Ocean, (f-j) Ross Sea, (k-o) Weddell Sea and (p-t) Adélie Land. Five columns represent five climate experiments (PI, mid-Holocene, Last Interglacial, LGM, and MIS3).   \label{sam_wmt}}
\end{figure}

\clearpage

\begin{figure}
\centering
\includegraphics[width=\textwidth]{figures/figS_sam_heat_density.pdf}
\caption{Austral winter (JJA) high-SAM minus low-SAM composites ($\pm$1.0$\sigma$ threshold) of the surface density flux produced by (a-e) the total heat flux, (f-j) its radiative component (shortwave plus longwave) and (k-o) its turbulent component (sensible plus latent heat). Five columns represent PI, the mid-Holocene, the Last Interglacial, LGM, and MIS3, respectively. The flux is converted with the thermal expansion coefficient at the winter surface temperature and salinity of each state, as in Fig. \ref{density_heat}, and shares one colour scale with Fig. \ref{sam_fwf}. Positive values mean that the surface water becomes denser. The dashed purple line in (a-e) marks where high-SAM winters lose 15~W~m$^{-2}$ more heat than low-SAM winters. Units: kg m$^{-2}$ month$^{-1}$.\label{sam_heat}}
\end{figure}

\clearpage

\begin{figure}
\centering
\includegraphics[width=\textwidth]{figures/figS_sam_fw_density.pdf}
\caption{Austral winter (JJA) high-SAM minus low-SAM composites ($\pm$1.0$\sigma$ threshold) of the surface density flux produced by (a-e) the total freshwater flux, (f-j) its sea ice component (brine rejection and melt) and (k-o) other sources (net precipitation plus runoff). Five columns represent PI, the mid-Holocene, the Last Interglacial, LGM, and MIS3, respectively. The flux is converted with the haline contraction coefficient and surface salinity of each state, as in Fig. \ref{density_fwf}, and shares one colour scale with Fig. \ref{sam_heat}. Positive values mean that the surface water becomes denser, that is, the ocean loses freshwater. Units: kg m$^{-2}$ month$^{-1}$.\label{sam_fwf}}
\end{figure}

\clearpage

\begin{figure}
\centering
\includegraphics[width=\textwidth]{figures/figR4_sam_push_pull.pdf}
\caption{The Southern Annular Mode acts on both the wind-driven upwelling and the sea ice salt pump, in every simulated climate state. All quantities are austral winter composites, high-SAM minus low-SAM years at the $\pm$1.2$\sigma$ threshold. (a) Anomaly of the zonally integrated northward Ekman transport at 60°S, computed from the simulated zonal wind stress. (b) Anomaly of the coastal brine-equivalent freshwater loss integrated south of 65°S, where a positive value denotes a stronger salt input to the ocean. (c) The resulting anomaly of the peak Southern Ocean transformation rate, with the density class of the maximum indicated. \label{sam_pushpull}}
\end{figure}

\clearpage

\begin{figure}
\centering
\includegraphics[width=\textwidth]{figures/figR5_sam_wind_ice.pdf}
\caption{Spatial expression of the Southern Annular Mode mechanism, shown as austral winter high-SAM minus low-SAM composites ($\pm$1.2$\sigma$ threshold) for the five climate states. (a, c, e, g, i) Zonal wind stress anomaly. (b, d, f, h, j) Sea ice concentration anomaly. Under the positive phase the westerlies strengthen over the circumpolar belt in every state. The sea ice concentration changes mainly near the winter ice edge, with gains and losses alternating between sectors and a net gain of ice area in every state. Near the coast the concentration hardly changes because the winter ice stays close to full cover there. The ice edge response lies further north and is about twice as large in the glacial states. \label{sam_windice}}
\end{figure}

\clearpage

\begin{figure}
      \centering\includegraphics[width=\textwidth]{figures/plot_wmt_4regions_5exps_annual_from_100years.pdf}
        \caption{Annual-mean water mass transformation rates [Sv] as a function of density $\sigma_2$ [kg m$^{-3}$] for (a-e) Southern Ocean, (f-j) Ross Sea, (k-o) Weddell Sea and (p-t) Adélie Land. Five columns represent five climate experiments (PI, mid-Holocene, Last Interglacial, LGM, and MIS3).  Lines show the mean values; shading denotes ±1  standard deviation across years.  \label{s2}}
    \end{figure}

\clearpage

\begin{figure}
\centering
\includegraphics[width=\textwidth]{figures/figR17_mld_pi_vs_obs.pdf}
\caption{Winter (JJA) mixed layer depth south of 60°S in (a) the pre-industrial simulation and (b) the ORAS5 ocean reanalysis \cite{Zuo2019}, averaged over 1990--2009. ORAS5 defines the mixed layer by a density threshold of 0.01~kg~m$^{-3}$ and the model by a bulk Richardson number criterion, so the point-by-point difference is not purely model error. The colour scale is finer between 0 and 200 m. Units: m. \label{mld_obs}}
\end{figure}

\begin{figure}
\centering
\includegraphics[width=\textwidth]{figures/figS_formation_location.pdf}
\caption{Where the simulated winter (JJA) surface transformation south of 60°S takes place. (a) Surface density flux that transforms water into the densest classes of each state, that is, the densities beyond which the transformation stays below 10\% of its maximum (threshold below each panel), on a logarithmic scale. (b) Surface density flux that transforms water across the density class of the transformation maximum. (c) Transformation in the class of the maximum, split into the continental shelf and slope (bottom depth shallower than 1000 m), deep water with a winter mixed layer deeper than 400 m, deep water under sea ice and ice-free deep water. (d) The same split for the densest classes. (e) The densest classes split by distance from the coast, with the part within 200 km of the coast divided at a winter sea ice concentration of 0.7. Lines in (a,b) mark the 1000 m isobath, the 400 m winter mixed layer depth and 60°S. Units: kg m$^{-2}$ month$^{-1}$ in (a,b) and Sv in (c-e). \label{formation_loc}}
\end{figure}

\end{appendices}

\clearpage
% ============================================================================
% HOW TO COMPILE:
% ============================================================================
% Run the following commands in order:
% 1. pdflatex manuscript.tex
% 2. bibtex manuscript
% 3. pdflatex manuscript.tex
% 4. pdflatex manuscript.tex
%
% Or use your LaTeX editor's "Build" button which usually does this automatically
% ============================================================================

\end{document}


